\section{Question and scope}

Can a translation-specialized language model provide a useful continuous input to a separate, frozen instruction-following language model?
The usual translation pipeline answers a Japanese question by first generating an English translation and then passing that text to an answering model.
Here, the alternative interface passes pooled hidden states from the translation model through a learned linear projection into the answering model's embedding space.
The source model does not generate an English translation on this path.
Both pretrained backbones remain fixed.

This interface could avoid committing to a single translated string, but that motivation does not establish an advantage.
A hidden-state sequence is neither a lossless encoding of the input nor an automatically compatible representation for another model.
Pooling discards information, and a causal decoder's early states cannot represent words that appear later in its input.
The projection must also learn an interface that the target model can use.
The empirical question is therefore whether this particular trained connection supports accurate free-running answers, and how its behavior compares with ordinary textual inputs.

The architectural precedent is LangBridge, which maps multilingual encoder states into a frozen language model through a learned interface~\citep{yoon2024langbridge}.
LangBridge studies cross-lingual transfer after alignment using English data.
The present study instead uses a frozen translation-specialized causal decoder, pooled intermediate states, and Japanese task-supervised fitting.
It is an exploratory adaptation of that general approach, rather than a reproduction of LangBridge or a new claim of zero-shot multilingual alignment.

The source is PLaMo 3 Translate 0.8B~\citep{pfn2026plamo}, and the target is Qwen2.5-1.5B-Instruct~\citep{qwen2024model,qwen2024report}.
Controlled Japanese arithmetic questions provide exact references and explicit train, validation and test partitions.
A separate template holdout tests sensitivity to surface wording.
The Japanese MGSM benchmark provides a distinct natural-language endpoint~\citep{shi2022mgsm}; its evaluation is reported separately from the generated tasks.
These endpoints address the behavior of a small, task-supervised system, not broad multilingual competence.

\section{Comparison logic}

The practical baselines give the frozen target a Japanese question, its paired English version, or an English translation generated by PLaMo.
The paired English version is a reference-language diagnostic rather than a deployable translation system.
The bridge receives task supervision that these frozen textual baselines do not.
Consequently, their comparison alone cannot isolate a benefit of continuous representations.
A learned soft-prefix baseline uses the same supervised training items while retaining the Japanese question as target-model text; it has fewer trainable parameters and is not parameter matched.

Layer and prefix-length ablations test specific interface choices.
An untrained projection tests the need for fitting the interface, while source-feature shuffling tests whether a trained system depends on the correspondence between each question and its prefix.
Neither control establishes where an arithmetic computation occurs.
In particular, source dependence does not by itself show that the target performs reasoning rather than using task-relevant information already present in the source states.
